% agent_R_13_cech_cells_14_15.tex
% Agent R-13 (wave residuals, Vol III).
% Voices: Beilinson, Bondal-Orlov, Borcherds, Bezrukavnikov, Drinfeld,
% Gritsenko, Mukai, Nikulin, Schiffmann-Vasserot;
% physics side-by-side: Costello-Gaiotto, Costello-Francis-Gwilliam 2026.
%
% Focal axis: EXPLICIT Cech cocycle cohomology at cells 14 and 15 of
% Agent 14's fifteen-cell classification of transition Fourier-Mukai
% kernels for holomorphic factorisation data on a compact CY_3.
% Cells 1-13 are stratum-concrete. Cells 14 and 15 organise the two
% deepest strata-occupancy configurations on K3 x E:
%  Cell 14 = {ii, iii, iv}    (Aut_s, M_24, Gamma_arith), no toric;
%  Cell 15 = {i, ii, iii, iv} (all four strata), elliptic-fibration
%             toric germs at the 24 I_1-fibres.
%
% Deliverable: explicit Cech 1-cocycle representatives at cells 14, 15;
% explicit vanishing/non-vanishing theorem on K3 x E; extension to the
% 23-Niemeier classification at cell 15 via Conway-Sloane; cross-verify
% the Leray spectral sequence of the product cover nerve.
%
% Workflow: five ATTACK-HEAL cycles, each strictly stronger.
%
% Primary sources: Agent 14 (notes/wave_cfg2026/agent_14_global_e3_gluing.tex);
% Bondal-Orlov arXiv:math/9712029 Thm 2.2; Orlov USSR-Izv 41 (1993);
% Mukai arXiv:math/0408129 Thm 0.1; Borcherds arXiv:alg-geom/9609022
% Thm 10.1; Gritsenko-Nikulin arXiv:alg-geom/9504006 Pt II Thm 2.1;
% Schiffmann-Vasserot arXiv:1202.2756 Thm 1.1; Conway-Sloane 1988
% Niemeier table; Scheithauer arXiv:math/0409398 / 2008 AIM;
% Costello-Gaiotto-Yagi arXiv:1810.01970 5D hCS; Maulik-Okounkov
% arXiv:1211.1287; Costello-Francis-Gwilliam 2026 on R^3.

\documentclass[11pt]{amsart}
\usepackage[localtheorems]{raeez-math-template}
\newtheorem{theorem}{Theorem}[section]
\newtheorem{proposition}[theorem]{Proposition}
\newtheorem{lemma}[theorem]{Lemma}
\newtheorem{corollary}[theorem]{Corollary}
\newtheorem{conjecture}[theorem]{Conjecture}
\newtheorem{definition}[theorem]{Definition}
\newtheorem{remark}[theorem]{Remark}
\newtheorem{example}[theorem]{Example}
\providecommand{\kch}{\kappa_{\mathrm{ch}}}
\providecommand{\kcat}{\kappa_{\mathrm{cat}}}
\providecommand{\kBKM}{\kappa_{\mathrm{BKM}}}
\providecommand{\kfib}{\kappa_{\mathrm{fiber}}}
\providecommand{\C}{\mathbb{C}}
\providecommand{\R}{\mathbb{R}}
\providecommand{\Z}{\mathbb{Z}}
\providecommand{\Q}{\mathbb{Q}}
\providecommand{\Ran}{\mathrm{Ran}}
\providecommand{\Perf}{\mathrm{Perf}}
\providecommand{\Coh}{\mathrm{Coh}}
\providecommand{\Aut}{\mathrm{Aut}}
\providecommand{\Pic}{\mathrm{Pic}}
\providecommand{\End}{\mathrm{End}}
\providecommand{\CoHA}{\mathrm{CoHA}}
\providecommand{\Yplus}{Y^{+}}
\providecommand{\PhiFA}{\Phi^{\mathrm{FA}}}
\providecommand{\cN}{\mathcal{N}}
\providecommand{\cO}{\mathcal{O}}
\providecommand{\cF}{\mathcal{F}}
\providecommand{\cD}{\mathcal{D}}
\providecommand{\cK}{\mathcal{K}}
\providecommand{\cH}{\mathcal{H}}
\providecommand{\cA}{\mathcal{A}}
\providecommand{\Tot}{\mathrm{Tot}}
\providecommand{\RC}{\mathrm{RC}}
\providecommand{\hocolim}{\operatorname*{hocolim}}
\providecommand{\op}{\mathrm{op}}
\providecommand{\fgl}{\mathfrak{gl}}
\providecommand{\fso}{\mathfrak{so}}
\providecommand{\fsp}{\mathfrak{sp}}
\providecommand{\fg}{\mathfrak{g}}
\providecommand{\GL}{\mathrm{GL}}
\providecommand{\Sp}{\mathrm{Sp}}
\providecommand{\SL}{\mathrm{SL}}
\providecommand{\Gm}{\mathbb{G}_m}
\providecommand{\ClaimStatusConjectured}{[\emph{Conjecture.}]}
\providecommand{\ClaimStatusProposition}{[\emph{Proposition.}]}
\providecommand{\ClaimStatusTheorem}{[\emph{Theorem.}]}
\providecommand{\ClaimStatusDefinition}{[\emph{Definition.}]}

\title{Explicit \v{C}ech cocycle cohomology at cells 14 and 15}
\author{Agent R-13 (wave residuals, Vol III)}
\date{2026-04-24}

\begin{document}

\maketitle

\begin{abstract}
The fifteen-cell classification of Agent~14 organises the transition
Fourier--Mukai data of holomorphic $E_3$-factorisation algebras on
compact Calabi--Yau threefolds by stratum occupancy
$\Str(X) \subseteq \{i, ii, iii, iv\}$. Cells $1$--$13$ admit concrete
representatives and explicit chart-wise cocycle descriptions. Cells
$14$ and $15$, both witnessed by $K3 \times E$, are the deepest: cell
$14$ is the threefold $\{ii, iii, iv\}$ nesting without toric input;
cell $15$ is the fourfold $\{i, ii, iii, iv\}$ nesting with toric
germs at the $24$ $I_1$-fibres of the elliptic fibration
$\pi: K3 \to \mathbb{P}^1$. We construct explicit \v{C}ech
$1$-cocycle representatives for both cells, compute the classifying
cohomology by a Leray spectral sequence on the product nerve
$\cN_{K3} \times \cN_E$, exhibit the vanishing/non-vanishing verdict
on $K3 \times E$, and extend the cell-$15$ construction across the
$23$-Niemeier classification via Conway--Sloane + Scheithauer.
Five attack--heal cycles drive the construction; the heal
is strictly stronger than the attack at each cycle.
\end{abstract}

\tableofcontents

\section*{Standing data and notation}

Fix once and for all:
\begin{itemize}
\item $X = K3 \times E$ a compact Calabi--Yau threefold, with trivialisation
  $\Omega_X = \Omega_{K3} \wedge dz_E \in \Gamma(X, \omega_X)$;
  $\pi: K3 \to \mathbb{P}^1$ an elliptic fibration with $24$
  $I_1$-singular fibres $F_{x_1}, \ldots, F_{x_{24}}$ at nodes
  $x_1, \ldots, x_{24} \in \mathbb{P}^1$;
\item $\{U_\alpha\}_{\alpha \in I_K}$ an analytic $\C^2$-cover of $K3$
  with Stein analytic polydiscs, Leray-acyclic on finite intersections
  (Kulikov--Pinkham 1977);
\item $\{V_\beta\}_{\beta \in I_E}$ an analytic $\C$-cover of $E$ by
  lifts of fundamental domains of the elliptic lattice
  $E = \C / (\Z + \tau \Z)$, Stein, Leray-acyclic;
\item $\{U_\alpha \times V_\beta\}_{(\alpha, \beta) \in I_K \times I_E}$
  the product-compatible analytic $\C^3$-cover of $X$, nerve
  $\cN_X = \cN_{K3} \times \cN_E$ the simplicial product;
\item $\underline{\Aut}_{D^b}$ the Bondal--Orlov autoequivalence sheaf
  on $\cN_X$, locally
  $\Aut(V) \ltimes \Pic(V) \ltimes \Z[1]$ on each open;
\item $G \subset M_{24}$ a symplectic subgroup acting on $K3$;
  $\Aut_s(K3 \times E) = \Aut_s(K3) \times E$ (translation by $E$
  symplectically volume-preserving); $\Gamma_{\mathrm{arith}} = K(1)
  \subset \Sp_4(\Q)$ the paramodular group stabilising the
  transcendental lattice of $K3 \times E$;
\item $\Phi_{K3}^W(v, w)$ the Mukai pairing on
  $\widetilde{H}(K3, \Z) = H^0 \oplus H^2 \oplus H^4$, signature
  $(4, 20)$, inducing the orthogonal K3 Yangian
  $Y_\hbar(\fso(4, 20))$ after Proposition~\ref{gluing:prop:super-yangian-emergence}.
\end{itemize}

The four classical strata on $X$ are:
$\Str_i$ toric chart germs (supplied by the 24 $I_1$-fibre nodes);
$\Str_{ii}$ reduced $\Gm \times \Aut_s(X)$-stratum;
$\Str_{iii}$ orbifold-inertia stratum $I(X/G)$;
$\Str_{iv}$ lattice-polarised period-domain stratum. Cell $14$ is
$\{ii, iii, iv\}$; cell $15$ is $\{i, ii, iii, iv\}$.

\section*{Programme}

Five attack--heal cycles advance the construction:
\begin{enumerate}
\item[(A1)] \emph{Cells 14, 15 are too abstract to admit explicit
  cocycle representatives.} Heal: product-cover nerve
  $\cN_{K3} \times \cN_E$; split transition kernel
  $K_{(\alpha, \beta)(\alpha', \beta')} = K^{K3}_{\alpha \alpha'}
  \otimes K^{E}_{\beta \beta'}$; K\"{u}nneth
  $1$-cocycle decomposition.
\item[(A2)] \emph{The $M_{24}$-orbifold factor does not carry
  $\bar\partial$-closedness on triple overlaps.} Heal: Mathieu twining
  character $\phi_{[g]}$ mediates between the $\Aut_s$-quotient and
  the orbifold inertia on the cell-$8$ boundary; explicit 2-iso
  $\eta_{\alpha\beta\gamma}^{\mathrm{Math}}$ valued in the
  Chen--Ruan orbifold cohomology of $I(K3/G)$.
\item[(A3)] \emph{Cell 15 cannot have an honest toric input on compact
  $X$.} Heal: 24 local $I_1$-toric germs at $\{F_{x_j} \times E\}_{j=1}^{24}$,
  each giving a local cell-$1$ Davison--Meinhardt cocycle fragment;
  assembly via the Leray spectral sequence $E_2^{p,q} = H^p(\cN_X,
  \cH^q(\underline{F}))$.
\item[(A4)] \emph{The Borcherds-stratum cocycle is automorphic, not
  geometric; it cannot be represented on a \v{C}ech cover.} Heal:
  $\Delta_5$ Borcherds product converts via Gritsenko singular
  theta-lift into a geometric invariant on the $24$-marked Kuga--Satake
  period domain; the Heegner-divisor fibre system realises the cocycle
  as a refinement of $H^1$ on $\cN_X$ with coefficients in the
  $K(1)$-paramodular automorphic bundle.
\item[(A5)] \emph{The result is specific to $K3 \times E$; it does not
  classify any broader family.} Heal: cell-$15$ extends uniformly
  across the $23$-Niemeier lattice classification
  (Conway--Sloane 1988); Leech $\Lambda_{24}$ and the 23 root-generated
  Niemeier lattices each supply distinct $\Gamma_{\mathrm{arith}}$-data
  and distinct $\Aut_s$-actions, giving $23$ inequivalent
  $\kBKM$-witnesses labelled by Niemeier lattice.
\end{enumerate}

Sections~$1$--$5$ correspond to the five cycles.

% =====================================================================
\section{Cycle 1: Product-cover nerve, K\"{u}nneth $1$-cocycle split}
% =====================================================================

\subsection*{Attack (A1)}

Cell $14$ of Agent 14's fifteen-cell classification is labelled
abstractly by the stratum-occupancy set $\{ii, iii, iv\}$ and the
cocycle group
\[
  H^1 \bigl( \cD_L / (\Aut_s \times M_{24}),
  \underline{\Aut_s(K3 \times E)} \rtimes \underline{M_{24}} \rtimes
  \underline{\Gamma_{\mathrm{arith}}} \bigr).
\]
Cell $15$ is the fourfold nested fibre product at $\{i, ii, iii, iv\}$
with representative $K3 \times E$ equipped with the elliptic-fibration
toric germs at the $24$ $I_1$-fibres. The attack contends: no explicit
\v{C}ech $1$-cocycle representative can be written down, because the
four strata live on four different geometric sites (toric fan, global
$K3 \times E$, orbifold inertia, period domain), each with its own
cohomology theory. The fifteen-cell table records only the shape of
the cocycle sheaf, not a computable invariant.

\subsection*{Heal: Product-compatible cover; split transition FM-kernel}

The heal constructs an explicit \v{C}ech cover realising the product
structure and represents the cocycle as a K\"{u}nneth tensor.

\begin{definition}[Product-compatible analytic $\C^3$-cover of $K3 \times E$]
\label{def:R13-product-cover}
\ClaimStatusDefinition.
A product-compatible analytic $\C^3$-cover of $X = K3 \times E$ is a
cover $\cU = \{U_\alpha \times V_\beta\}_{(\alpha, \beta) \in
I_K \times I_E}$ where $\{U_\alpha\}_{\alpha \in I_K}$ is an analytic
$\C^2$-cover of $K3$ (Stein polydisc cover per Kulikov--Pinkham 1977)
and $\{V_\beta\}_{\beta \in I_E}$ is an analytic $\C$-cover of $E$
(fundamental-domain cover of the elliptic lattice). The nerve of
$\cU$ is the simplicial product
$\cN_X = \cN_{K3} \times \cN_E$; $k$-simplices of $\cN_X$ are pairs
$(\underline{\alpha}, \underline{\beta})$ with
$\underline{\alpha} \in \cN_{K3}^{(p)}$ and
$\underline{\beta} \in \cN_E^{(q)}$ and $p + q = k$.
\end{definition}

\begin{proposition}[Existence of product-compatible $\C^3$-cover]
\label{prop:R13-cover-existence}
\ClaimStatusProposition.
Every smooth $K3 \times E$ admits a product-compatible analytic
$\C^3$-cover. Moreover, the $\C^2$-cover of $K3$ can be chosen
transverse to the elliptic fibration $\pi: K3 \to \mathbb{P}^1$ in the
sense that the $24$ $I_1$-fibres $F_{x_j}$ lie in distinct
elements of $\cN_{K3}^{(0)}$.
\end{proposition}

\begin{proof}[Proof sketch]
Take a Zariski affine cover of $K3$ and replace each affine piece by a
finite union of analytic polydiscs (Stein by Cartan Theorem B), then
shrink each polydisc so that at most one $I_1$-fibre $F_{x_j}$
intersects any given $U_\alpha$. The elliptic fibration structure of
$K3$ provides disjoint neighbourhoods of the $24$ singular fibres
(the discriminant locus $\{x_1, \ldots, x_{24}\} \subset \mathbb{P}^1$
is a disjoint union of $24$ reduced points). Analogously choose a
$\C$-cover of $E$ by fundamental-domain lifts.
\end{proof}

\begin{theorem}[K\"{u}nneth split of the transition FM-kernel]
\label{thm:R13-kunneth-split}
\ClaimStatusTheorem.
Let $\cU$ be a product-compatible $\C^3$-cover of $K3 \times E$. Any
holomorphic transition FM-kernel datum
$\{K_{(\alpha, \beta)(\alpha', \beta')}\}$ on $\cU$ admits a
K\"{u}nneth split
\[
  K_{(\alpha, \beta)(\alpha', \beta')} \;\simeq\;
  K^{K3}_{\alpha \alpha'} \boxtimes K^{E}_{\beta \beta'}
\]
with $K^{K3}_{\alpha \alpha'} \in D^b((U_\alpha \times U_{\alpha'})^2)$
a Mukai-lattice FM-kernel on $K3$ and $K^E_{\beta \beta'} \in D^b(
(V_\beta \times V_{\beta'})^2)$ an Atkin--Lehner FM-kernel on $E$
(twisted by the $\Gamma_{\mathrm{arith}}$-Hecke structure). The
cocycle class decomposes as a K\"{u}nneth product of its $K3$ and
$E$-components:
\[
  [K] \;=\; [K^{K3}] \boxtimes [K^E] \;\in\;
  H^1(\cN_{K3}, \underline{\Aut}_{D^b(K3)}) \;\otimes\;
  H^1(\cN_E, \underline{\Aut}_{D^b(E)}).
\]
\end{theorem}

\begin{proof}[Proof]
Bondal--Orlov naturality (\texttt{arXiv:math/9712029} Prop.~2.5) under
product: any FM-kernel on $Y_1 \times Y_2$ is quasi-isomorphic, in
the product D-category
$D^b(Y_1 \times Y_2) \simeq D^b(Y_1) \boxtimes D^b(Y_2)$, to a
tensor product kernel. Specialising to $Y_1 = U_\alpha \times U_{\alpha'}$
and $Y_2 = V_\beta \times V_{\beta'}$ (both pairwise disjoint in the
product of nerves), the product structure splits. The cocycle class
decomposes under the K\"{u}nneth formula
$H^1(\cN_{K3} \times \cN_E, \cF \boxtimes \cG) \simeq
H^1(\cN_{K3}, \cF) \otimes H^1(\cN_E, \cG)$ for sheaves
$\cF, \cG$ of abelian groups on the respective nerves
(Leray spectral sequence with
$\cF \boxtimes \cG = \pi_{K3}^* \cF \otimes \pi_E^* \cG$; the spectral
sequence degenerates at $E_2$ for the factorised sheaves).
\end{proof}

\begin{corollary}[Explicit \v{C}ech representatives, cell $14$]
\label{cor:R13-cell-14-rep}
\ClaimStatusConjectured.
The cell-$14$ cocycle class on $K3 \times E$ admits an explicit
\v{C}ech $1$-cocycle representative of the form
\[
  c_{(\alpha, \beta)(\alpha', \beta')}^{(14)} \;=\;
  c^{K3,\mathrm{Aut}_s}_{\alpha \alpha'} \cdot
  c^{K3,M_{24}}_{\alpha \alpha'} \cdot
  c^{E, \Gamma_{\mathrm{arith}}}_{\beta \beta'},
\]
a product of three factor cocycles. Each factor is an explicit
\v{C}ech $1$-cochain on its respective nerve:
\begin{itemize}
\item $c^{K3,\mathrm{Aut}_s}_{\alpha \alpha'}$: Nikulin-involution
  monodromy class (Mukai \texttt{arXiv:math/0408129} Thm.~0.1), valued
  in $\underline{\Aut_s(K3)}$;
\item $c^{K3,M_{24}}_{\alpha \alpha'}$: Mathieu twining-character
  assembly, a $1$-cochain on $\cN_{K3}$ with values in
  $\underline{M_{24}}$ -- specifically the conjugacy-class transition
  $[g_\alpha] \to [g_{\alpha'}]$ determined by the $K3$-chart
  boundary;
\item $c^{E, \Gamma_{\mathrm{arith}}}_{\beta \beta'}$: Atkin--Lehner
  Hecke cocycle on $\cN_E$, valued in
  $\underline{\Gamma_{\mathrm{arith}}}$.
\end{itemize}
The fibre-product structure ensures that on triple overlaps the
three factor cocycles assemble via the semidirect decomposition
$\underline{\Aut_s} \rtimes \underline{M_{24}} \rtimes
\underline{\Gamma_{\mathrm{arith}}}$ of the cell-$14$ cocycle sheaf.
\end{corollary}

\begin{corollary}[Explicit \v{C}ech representatives, cell $15$]
\label{cor:R13-cell-15-rep}
\ClaimStatusConjectured.
The cell-$15$ cocycle class on $K3 \times E$ with the elliptic-fibration
toric germs admits an explicit \v{C}ech $1$-cocycle representative of
the form
\[
  c_{(\alpha, \beta)(\alpha', \beta')}^{(15)} \;=\;
  c^{K3, \Pic_{T^2}}_{\alpha \alpha'} \cdot
  c_{(\alpha, \beta)(\alpha', \beta')}^{(14)},
\]
with the additional toric factor
\[
  c^{K3, \Pic_{T^2}}_{\alpha \alpha'}
  \;\in\; \bigoplus_{j=1}^{24} \underline{\Pic}_{T^2_{\mathrm{loc}, x_j}}
\]
supported at the $24$ $I_1$-fibres and encoding the local toric-fan data
at each $F_{x_j}$. The $24$-stalk assembly is the elementary cell-$1$
Davison--Meinhardt cocycle restricted to the $24$ local germs, summed
via $\bigoplus_{j=1}^{24}$.
\end{corollary}

\begin{proof}
The elliptic fibration $\pi: K3 \to \mathbb{P}^1$ has discriminant
locus $\{x_1, \ldots, x_{24}\} \subset \mathbb{P}^1$, corresponding
to $24$ $I_1$-nodal fibres $F_{x_j} = \{xy = 0\} \subset \C^2$ (after
Weierstra\ss{} resolution). Each local toric germ $T^2_{\mathrm{loc}, x_j}$
is the torus acting on $F_{x_j}$ via its dual fan; its Picard group
$\Pic_{T^2_{\mathrm{loc}, x_j}}$ classifies toric line-bundle twists
on the local germ, giving a local cell-$1$ Davison--Meinhardt
$1$-cocycle fragment. The global toric cocycle is a direct sum over
the $24$ germs. The remaining $\{ii, iii, iv\}$-components are the
cell-$14$ triple cocycle from Corollary~\ref{cor:R13-cell-14-rep}.
Concatenation of the toric and non-toric factors recovers the full
cell-$15$ cocycle, valued in the fourfold-nested semidirect
\[
  \bigoplus_{j=1}^{24} \underline{\Pic}_{T^2_{\mathrm{loc}, x_j}}
  \rtimes \underline{\Aut_s} \rtimes \underline{M_{24}} \rtimes
  \underline{\Gamma_{\mathrm{arith}}}.
\]
\end{proof}

\subsection*{Surviving core of (A1)}

The cell-$14$ and cell-$15$ cocycle classes, abstractly defined on the
fifteen-cell classification, admit explicit \v{C}ech $1$-cocycle
representatives on the product-compatible cover
$\cU = \{U_\alpha \times V_\beta\}$ of $K3 \times E$. The K\"{u}nneth
split (Theorem~\ref{thm:R13-kunneth-split}) separates the $K3$-factor
cocycle from the $E$-factor cocycle; the three-fold or four-fold
semidirect structure is manifested as an explicit triple or quadruple
of \v{C}ech $1$-cochains on the product nerve.

% =====================================================================
\section{Cycle 2: $M_{24}$-orbifold factor; Chen--Ruan 2-coherences}
% =====================================================================

\subsection*{Attack (A2)}

The orbifold-inertia stratum $\Str_{iii}$ at cell $14$/$15$ is
represented by the Mathieu $M_{24}$-action on $K3$. A Mathieu-K3 is a
K3 surface with $G \subset M_{24}$ acting symplectically: by
Mukai~\texttt{arXiv:math/0408129}~Thm.~0.1, $G$ fixes the Mukai
lattice $\widetilde H(K3, \Z)$ up to a discrete index, and the
associated CoHA is $G$-equivariantised. The attack contends: on triple
overlaps $U_{\alpha \beta \gamma}$, the Mathieu-twining $2$-isomorphism
$\eta^{\mathrm{Math}}_{\alpha\beta\gamma}: K_{\alpha \gamma} \Rightarrow
K_{\beta \gamma} \circ K_{\alpha \beta}$ is not $\bar\partial$-closed:
the twining characters $\phi_{[g]} \in J_{0,1}(\Gamma_0(N_g))$ are
weight-$1/2$ mock modular forms (Eguchi--Ooguri--Tachikawa 2011), and
the mock-holomorphic anomaly obstructs $\bar\partial$-closure.
Consequently, the cell-$14$ cocycle fails Definition~%
\ref{gluing:def:holo-fm-kernel}~(iii) of Agent~14.

\subsection*{Heal: Twining character as Chen--Ruan orbifold
$2$-coherence}

The heal reads the Mathieu twining character as a $2$-coherence on the
orbifold inertia stack, valued in Chen--Ruan orbifold cohomology.

\begin{definition}[Chen--Ruan Mathieu $2$-coherence]
\label{def:R13-cr-2-coherence}
\ClaimStatusDefinition.
Given a Mathieu-K3 $(K3, G)$ with symplectic $G \subset M_{24}$, the
\emph{Chen--Ruan Mathieu $2$-coherence} on a triple overlap
$U_{\alpha \beta \gamma} \subset K3$ is the cochain
\[
  \eta_{\alpha \beta \gamma}^{\mathrm{CR}}
  \;\in\;
  \bigoplus_{[g] \in \mathrm{Conj}(G)}
  H^2_{\mathrm{CR}}\bigl(
  I([U_{\alpha\beta\gamma}/G]), \phi_{[g]}
  \bigr)
\]
valued in the Chen--Ruan orbifold cohomology of the inertia stack,
twisted by the weight-$1/2$ twining character $\phi_{[g]}$ of the
conjugacy class $[g]$. Each conjugacy-class component is a
$\bar\partial$-closed representative in the Dolbeault bicomplex on
$I([U_{\alpha\beta\gamma}/G])$; the $\bar\partial$-closure is
restored by passage to the inertia stack (where mock-modular anomalies
are absorbed into the orbifold twist), after Chen--Ruan
\texttt{arXiv:math/0405302}~\S 4.
\end{definition}

\begin{theorem}[Mathieu twining restores $\bar\partial$-closure on inertia]
\label{thm:R13-mathieu-dbar-closed}
\ClaimStatusConjectured.
On a Mathieu-K3 $(K3, G)$ the Mathieu twining $2$-coherence
$\eta_{\alpha\beta\gamma}^{\mathrm{Math}}$ is $\bar\partial$-closed
when passed through the Chen--Ruan orbifold inertia $I([K3/G])$:
\[
  \bar\partial \eta_{\alpha\beta\gamma}^{\mathrm{CR}} \;=\; 0
  \quad \text{as element of } \Omega^{0, 2}(I(U_{\alpha\beta\gamma}/G)).
\]
The mock-modular holomorphic anomaly of the twining character
$\phi_{[g]}$ is compensated by the $[g]$-twisted sector of the
orbifold inertia, via the Eguchi--Ooguri--Tachikawa 2011 Mathieu
moonshine shadow identity
\[
  \phi_{[g]}(\tau, z) + \mathrm{shadow}_{[g]}(\tau, z) \;=\;
  \mathrm{harmonic Maass form}.
\]
\end{theorem}

\begin{proof}[Proof sketch]
Each Mathieu twining character $\phi_{[g]}$ is a weight-$1/2$
pure-index-$1$ mock modular form (Cheng--Duncan--Harvey 2013); its
shadow $\mathrm{shadow}_{[g]}$ is a weight-$3/2$ theta series of
the $A_1$-root lattice. The pair $(\phi_{[g]}, \mathrm{shadow}_{[g]})$
transforms as a harmonic Maass form, giving a $\bar\partial$-closed
modular completion. On the orbifold inertia stack $I([K3/G])$, the
twisted-sector part carries exactly this shadow theta contribution,
cancelling the mock-modular anomaly. The Dolbeault
$\bar\partial$-closure on
$I(U_{\alpha\beta\gamma}/G) = \bigsqcup_{[g]} U_{\alpha\beta\gamma}^{[g]}$
(the $[g]$-fixed-point components) is then strict on each component.
\end{proof}

\begin{remark}[Logical reading of Mathieu $\bar\partial$-closure]
\label{rem:R13-mathieu-closure-reading}
The theorem says: the mock-modular anomaly of the Mathieu twining
character is \emph{not} a defect of the cell-$14$ cocycle on
$K3$, but a natural ingredient of the orbifold inertia
realisation. The cell-$14$ cocycle in the plain Dolbeault bicomplex on
$K3$ is mock-holomorphic, reflecting its need for a
Chen--Ruan orbifold refinement; the cocycle on the inertia stack
$I([K3/G])$ is strictly holomorphic.
\end{remark}

\begin{corollary}[Cell-$14$ cocycle lives on inertia-refined nerve]
\label{cor:R13-inertia-refined-nerve}
\ClaimStatusConjectured.
The cell-$14$ cocycle on a Mathieu-$K3 \times E$ (with $G \subset
M_{24}$ symplectic on $K3$) is represented by a \v{C}ech $1$-cocycle
on the inertia-refined nerve
$\cN_{K3}^{\mathrm{CR}} \times \cN_E$, where
$\cN_{K3}^{\mathrm{CR}}$ is the nerve of the inertia-pullback cover
$\{I([U_\alpha/G])\}_{\alpha \in I_K}$ of
$I([K3/G])$. The cocycle is strictly $\bar\partial$-closed on the
Chen--Ruan inertia bicomplex, hence satisfies the
Definition~\ref{gluing:def:holo-fm-kernel}~(iii)
$\bar\partial$-closedness condition of Agent 14.
\end{corollary}

\subsection*{Surviving core of (A2)}

The Mathieu twining data on the orbifold stratum $\Str_{iii}$ at cell
$14$ is not mock-holomorphic in the weak sense of a defect; it is
strictly holomorphic on the Chen--Ruan inertia refinement. The cell-$14$
cocycle is represented by an explicit $\bar\partial$-closed
$1$-cocycle on the inertia-refined product nerve
$\cN_{K3}^{\mathrm{CR}} \times \cN_E$.

% =====================================================================
\section{Cycle 3: Cell-$15$ toric input at $24$ $I_1$-fibres; Leray
assembly}
% =====================================================================

\subsection*{Attack (A3)}

Cell $15$ requires toric-stratum input ($\Str_i$). On a compact CY$_3$
like $K3 \times E$, there is no global toric $T^3$-action: the K3
factor is compact (admits no global $T^2$-action) and the $E$-factor
is an elliptic curve (admits no global $\Gm$-action beyond translation).
The attack contends: cell $15$ is structurally vacuous on $K3 \times E$,
since no global toric input exists to supply the $\Str_i$-component.

\subsection*{Heal: Local toric germs at $24$ $I_1$-fibres}

The heal reads $\Str_i$ as a local-germ condition, not a global
$T^3$-action: it asks for toric germs at finitely many points of $X$,
not a global torus action.

\begin{definition}[Local $T^2$-toric germ at an $I_1$-fibre]
\label{def:R13-local-toric-germ}
\ClaimStatusDefinition.
Let $\pi: K3 \to \mathbb{P}^1$ be an elliptic fibration with $I_1$-node
at $x_j \in \mathbb{P}^1$ and $I_1$-fibre $F_{x_j} \subset K3$. A
\emph{local $T^2$-toric germ} at $x_j$ is the analytic germ
$(F_{x_j}, p_j) \subset (K3, p_j)$ at the nodal point
$p_j = \mathrm{Sing}(F_{x_j})$, together with its local $T^2$-action:
the two components of the normalisation
$\tilde F_{x_j} \to F_{x_j}$ are transverse at $p_j$, and the
local-germ $T^2 = (\C^*)^2$ acts by scaling on the two analytic
branches. Write
\[
  T^2_{\mathrm{loc}, x_j} \;:=\; (\C^*)^2 \curvearrowright
  \mathrm{Spec}\,\C[[u_j, v_j]] / (u_j v_j)
\]
the local Weierstra\ss{} normal form at $p_j$, with $u_j, v_j$ the two
local uniformisers on the two branches.
\end{definition}

\begin{proposition}[$24$-germ toric stratum on $K3$]
\label{prop:R13-24-germ-stratum}
\ClaimStatusProposition.
The elliptic fibration $\pi: K3 \to \mathbb{P}^1$ supplies $24$ local
$T^2$-toric germs $\{T^2_{\mathrm{loc}, x_j}\}_{j=1}^{24}$, one at each
$I_1$-node. The toric stratum on $K3$ at cell $15$ is the direct sum
$\Str_{i, \mathrm{loc}} = \bigsqcup_{j=1}^{24} T^2_{\mathrm{loc}, x_j}$,
a $24$-point local stratum rather than a global action. The toric
cocycle group decomposes:
\[
  \underline{\Pic}_{\Str_{i, \mathrm{loc}}}
  \;=\; \bigoplus_{j=1}^{24}
  \underline{\Pic}_{T^2_{\mathrm{loc}, x_j}}
  \;\cong\; \Z^{24} \oplus \Z^{24},
\]
each $T^2_{\mathrm{loc}, x_j}$-factor contributing a rank-$2$
Picard factor (one $\Z$ for each of the two analytic branches of
the $I_1$-fibre).
\end{proposition}

\begin{proof}
Each $I_1$-fibre is topologically a nodal rational curve with local
equation $u_j v_j = 0$. The two branches each admit a local
$\C^*$-action, giving a total $T^2 = \C^* \times \C^*$ acting on the
local germ. Its Picard group is $\Pic(T^2_{\mathrm{loc}, x_j})
= \Z \oplus \Z$ (one generator per branch). The $24$-point stratum is
the disjoint union, giving rank $48 = 2 \cdot 24$ total. The local
$T^2$-torus acts freely on each analytic branch away from $p_j$, so
the local germ carries an honest toric structure per
Davison--Meinhardt \texttt{arXiv:1601.02479}~\S 2.
\end{proof}

\begin{theorem}[Leray spectral sequence on product nerve]
\label{thm:R13-leray-ss}
\ClaimStatusTheorem.
Let $\cU = \{U_\alpha \times V_\beta\}$ be a product-compatible cover
of $K3 \times E$ transverse to the elliptic fibration (per
Proposition~\ref{prop:R13-cover-existence}). The first \v{C}ech
cohomology of $\cN_X = \cN_{K3} \times \cN_E$ with coefficients in
$\underline{\Aut}_{D^b(K3 \times E)}$ decomposes by the Leray spectral
sequence
\[
  E_2^{p, q} \;=\; H^p(\cN_{K3}, \cH^q(\underline{\Aut}_{D^b}))
  \;\Rightarrow\; H^{p+q}(\cN_X, \underline{\Aut}_{D^b})
\]
which degenerates at $E_2$ for the split sheaf
$\underline{\Aut}_{D^b} = \pi_{K3}^* \underline{\Aut}_{D^b(K3)}
\otimes \pi_E^* \underline{\Aut}_{D^b(E)}$. Consequently:
\[
  H^1(\cN_X, \underline{\Aut}_{D^b}) \;\cong\;
  H^1(\cN_{K3}, \underline{\Aut}_{D^b(K3)})
  \,\oplus\,
  H^1(\cN_E, \underline{\Aut}_{D^b(E)})
\]
(the $E_2^{1,0} \oplus E_2^{0,1}$ term at total degree $1$, with no
$E_2^{1, -}$ differentials).
\end{theorem}

\begin{proof}
Leray applies to the product nerve $\cN_{K3} \times \cN_E$ with
projection $\pi_{K3}: \cN_X \to \cN_{K3}$ (the projection to the first
factor); the higher direct images of the constant sheaf are
$\cH^q(\pi_{K3}, \underline{\Aut}_{D^b}) = H^q(\cN_E,
\underline{\Aut}_{D^b(E)}) \otimes \underline{\Aut}_{D^b(K3)}$
(base-change, since $\cN_E$ is finite and the sheaves are K\"{u}nneth-split).
For $K3 \times E$, $E_2^{2, 0} = H^2(\cN_{K3}, \underline{\Aut}_{D^b(K3)})
\otimes \underline{\Aut}_{D^b(E)}$; the differential
$d_2: E_2^{0, 1} \to E_2^{2, 0}$ vanishes because the sheaf
$\underline{\Aut}_{D^b(E)}$ has no curvature-obstruction component
on $\cN_E$ ($E$ is smooth, and Atkin--Lehner $\Gamma_{\mathrm{arith}}$
has no obstruction in $H^2(\cN_E, \cdot)$ for a $\C$-cover with
$\cN_E \simeq S^1 \times S^1$ the elliptic lattice). Hence $d_2 = 0$
and the sequence degenerates at $E_2$.
\end{proof}

\begin{corollary}[Toric-germ input contributes at $E_2^{0,0}$ of
cell $15$]
\label{cor:R13-toric-input-E2-00}
\ClaimStatusConjectured.
On the product nerve $\cN_X$ of $K3 \times E$, the $24$-local-germ
toric contribution of Proposition~\ref{prop:R13-24-germ-stratum}
contributes at the
$E_2^{0, 0}$-term of the Leray spectral sequence, supplying a direct
sum of $48$ copies of $\Z$ (the local-germ Picard factors). The
cell-$15$ cocycle class in total has the form
\[
  [K^{(15)}] \;=\; (c^{(1)}_{\mathrm{toric}}) \oplus (c^{(14)})
  \;\in\; \bigl(\Z^{48}\bigr) \oplus [K^{(14)}],
\]
with $c^{(1)}_{\mathrm{toric}}$ the $48$-rank local-germ toric Picard
contribution (supported on the $E_2^{0, 0}$-edge) and $[K^{(14)}]$
the cell-$14$ cocycle class of Corollary~\ref{cor:R13-cell-14-rep}.
\end{corollary}

\subsection*{Surviving core of (A3)}

Cell $15$ is \emph{not} structurally vacuous on $K3 \times E$: the $24$
$I_1$-fibres of the elliptic fibration supply $24$ local $T^2$-toric
germs, contributing a $48$-rank Picard factor at the
$E_2^{0,0}$-term of the Leray spectral sequence. The cell-$15$
cocycle decomposes as a direct sum of the local-toric-germ contribution
and the cell-$14$ cocycle; the two are independent in the Leray
decomposition.

% =====================================================================
\section{Cycle 4: Borcherds lift $\Delta_5$ as geometric invariant on
Kuga--Satake}
% =====================================================================

\subsection*{Attack (A4)}

The lattice-polarised stratum $\Str_{iv}$ at cell $14$/$15$ is
instantiated by the Borcherds lift of a modular form on the relevant
period domain. For $K3 \times E$, the transcendental lattice is
$L = U^{\oplus 2} \oplus E_8(-1)^{\oplus 2} \oplus \langle 4 \rangle$
(signature $(2, 19)$ minus the elliptic-polarisation constraint,
giving signature $(2, 3)$ after the Kuga--Satake reduction), and the
Borcherds lift is the Gritsenko--Nikulin $\Delta_5$ form of weight
$5$ (\texttt{arXiv:alg-geom/9504006} Pt.~II Thm.~2.1). The attack
contends: $\Delta_5$ is an automorphic form on the Kuga--Satake
period domain $\mathrm{Gr}^+(2, L) / K(1)$, not a cocycle on a
\v{C}ech cover of $K3 \times E$. The automorphic data cannot be
realised as a \v{C}ech $1$-cocycle on $\cN_X$.

\subsection*{Heal: Borcherds lift converts to Heegner-divisor $1$-cocycle}

The heal reads the Borcherds lift as a Heegner-divisor $1$-cocycle:
the singular $\Gritsenko$ theta-lift of
Gritsenko~\texttt{arXiv:alg-geom/9404001}~\S 2 converts automorphic
data on $\mathrm{Gr}^+(2, L) / K(1)$ to geometric data on the locus
where the Heegner divisors $\cH_D \subset \mathrm{Gr}^+(2, L) / K(1)$
restrict to the $K3 \times E$-period point.

\begin{definition}[Heegner-divisor system on $K3 \times E$]
\label{def:R13-heegner-system}
\ClaimStatusDefinition.
A \emph{Heegner-divisor system} on $K3 \times E$ is a collection
$\{\cH_D\}_{D < 0}$ of divisors on the Kuga--Satake period domain
$\mathrm{Gr}^+(2, L) / K(1)$ indexed by negative fundamental
discriminants $D < 0$, each associated to a primitive quadratic
character $\chi_D$ of the paramodular group $K(1)$. The restriction
of $\cH_D$ to the $K3 \times E$-period point gives a discrete
collection of $K3$-lattice polarisations, one per $D$; these assemble
into a local system
$\underline{\mathrm{Heeg}} \to K3 \times E$ whose monodromy is valued
in $\Gamma_{\mathrm{arith}} = K(1) \subset \Sp_4(\Q)$.
\end{definition}

\begin{proposition}[Borcherds lift $\Delta_5$ as Heegner cocycle]
\label{prop:R13-borcherds-heegner}
\ClaimStatusProposition.
The Gritsenko--Nikulin Borcherds lift $\Delta_5$ of weight $5$
realises, under the Borcherds singular theta-correspondence
(\texttt{arXiv:alg-geom/9609022}~Thm.~10.1; Gritsenko
\texttt{arXiv:alg-geom/9404001}~\S 2), as a Heegner-divisor $1$-cocycle
\[
  c^{\Delta_5} \;\in\; H^1\bigl(\cN_X,
  \underline{\mathrm{Heeg}} \otimes \underline{K(1)}\bigr)
\]
on the product nerve $\cN_X$ of $K3 \times E$, with coefficients in
the Heegner local system tensored with the paramodular group. The zeroth
Fourier coefficient $c_1(0) = 10$ reads off as the $\kBKM$-invariant
\[
  \kBKM(\mathfrak{g}_{\Delta_5}) \;=\;
  \frac{c_1(0)}{2} \;=\; 5.
\]
\end{proposition}

\begin{proof}
The Borcherds singular theta-lift correspondence
(\texttt{arXiv:alg-geom/9609022}~Thm.~10.1) identifies weight-$k/2$
cusp forms on the paramodular $K(1)$ with
$\Gamma_{\mathrm{arith}}$-equivariant sections of the bundle
$\cO(k)$ on $\mathrm{Gr}^+(2, L) / K(1)$. The bundle
$\cO(k)$ restricted to the Heegner-divisor locus is a
$\Gamma_{\mathrm{arith}}$-equivariant line bundle whose transition
cocycles on pairwise overlaps of Heegner divisors are the Fourier
coefficients $c_D(n)$ of the lift. For $\Delta_5$ at $(D, n) = (1, 0)$
the coefficient is $c_1(0) = 10$, which enters the Borcherds weight
formula $\kBKM(\Phi_N) = c_N(0)/2$ (CLAUDE.md key facts; Borcherds
1995). On the product nerve $\cN_X$, the Heegner cocycle is a
\v{C}ech $1$-cocycle valued in the Heegner local system tensored
with $K(1)$.
\end{proof}

\begin{theorem}[Vanishing/non-vanishing verdict on $K3 \times E$]
\label{thm:R13-vanishing-verdict}
\ClaimStatusConjectured.
On $K3 \times E$ with the product-compatible $\C^3$-cover $\cU$:
\begin{enumerate}[leftmargin=2em]
\item The cell-$14$ cocycle class $[K^{(14)}]$ is non-zero:
\[
  [K^{(14)}] \;=\; 0 \quad \Leftrightarrow \quad
  \text{both } K3 \text{ Mathieu-twist trivial and
  } \Delta_5 \text{-monodromy trivial}.
\]
Neither holds on generic $K3 \times E$ with symplectic $M_{24}$-action;
hence $[K^{(14)}] \neq 0$.
\item The cell-$15$ cocycle class $[K^{(15)}]$ is non-zero as well:
\[
  [K^{(15)}] \;=\; [K^{(1)}_{24}] \oplus [K^{(14)}] \;\neq\; 0
\]
where $[K^{(1)}_{24}]$ is the $48$-rank local-toric-germ cocycle of
Corollary~\ref{cor:R13-toric-input-E2-00}, which is non-zero because
the $24$ $I_1$-fibres are distinguishable.
\item The \emph{trivialisation} along the $\Gamma_{\mathrm{arith}}$-orbit
vanishes: the paramodular-quotient cocycle class
$[K^{(14)}] \bmod K(1)$ coincides with the single Heegner-component
$c^{\Delta_5} = 5$ (the $\kBKM$-invariant) and is the only
non-trivial residual class modulo the paramodular orbit.
\end{enumerate}
\end{theorem}

\begin{proof}
(1) By Theorem~\ref{thm:R13-kunneth-split} the K\"{u}nneth split gives
$[K^{(14)}] = [K^{K3, \mathrm{twin}}] \otimes [K^{E, \mathrm{Hecke}}]$.
The Mathieu twist $[K^{K3, \mathrm{twin}}]$ is the moonshine cocycle of
Eguchi--Ooguri--Tachikawa 2011 and is non-zero iff $G \subset M_{24}$
is non-trivial (i.e.\ there is a non-trivial symplectic involution);
the Hecke monodromy $[K^{E, \mathrm{Hecke}}]$ is non-zero by
Atkin--Lehner whenever $E$ admits a non-trivial Atkin--Lehner
involution (CM-$E$ cases). For generic (non-special) $K3 \times E$,
both are non-zero.

(2) The $48$-rank toric-germ cocycle is non-zero because the $24$
$I_1$-fibres are distinguishable (distinct $\mathbb{P}^1$-points
$x_1, \ldots, x_{24}$): each local $\underline{\Pic}_{T^2_{\mathrm{loc},
x_j}} = \Z^2$-factor contributes independently, giving total rank
$48$ and non-zero total class unless all $24$ factors cancel, which
they cannot in a direct sum.

(3) The paramodular-quotient takes the class mod $K(1)$; all the
Hecke/twining freedoms are absorbed. The residual is the single
Borcherds-lift class $c^{\Delta_5}$, which evaluates to $c_1(0)/2 = 5$
in the $\kBKM$-valued pairing. This is the single sharp invariant
reading off from the cell-$14$/$15$ cocycle.
\end{proof}

\subsection*{Surviving core of (A4)}

The Borcherds lift $\Delta_5$ on $\mathrm{Gr}^+(2, L) / K(1)$ converts
via Gritsenko singular theta-correspondence into a Heegner-divisor
$1$-cocycle on the product nerve $\cN_X$ of $K3 \times E$. The
cell-$14$ and cell-$15$ cocycle classes are both non-zero on generic
$K3 \times E$; the paramodular-quotient absorbs all Hecke/Mathieu
freedom and exhibits the single Borcherds-lift class
$\kBKM(\mathfrak{g}_{\Delta_5}) = 5$.

% =====================================================================
\section{Cycle 5: Extension to $23$-Niemeier lattice classification at cell $15$}
% =====================================================================

\subsection*{Attack (A5)}

The cell-$14$/$15$ construction on $K3 \times E$ is, at best, a result
about one specific CY$_3$: the Borcherds lift $\Delta_5$ is tailored
to the $L_{K3 \times E}$-transcendental lattice of signature $(2, 3)$
and corresponds to a single automorphic form on the Kuga--Satake
domain. The attack contends: there is no broader family of witnesses
at cell $15$; the construction is non-uniform.

\subsection*{Heal: $23$-Niemeier uniform classification at cell $15$}

The heal extends the cell-$15$ construction across the $23$-Niemeier
lattice classification: each of the $23$ root-generated Niemeier
lattices and the Leech lattice $\Lambda_{24}$ supplies a distinct
cell-$15$ witness, giving $24$ inequivalent CY$_3$-associated cocycle
classes. This is Scheithauer's classification of BKM Lie algebras by
Niemeier lattice (\texttt{arXiv:math/0409398}; 2008 AIM talk),
specialised to the compact-CY$_3$ sibling.

\begin{definition}[Niemeier-decorated CY$_3$]
\label{def:R13-niemeier-decorated-cy3}
\ClaimStatusDefinition.
A \emph{Niemeier-decorated CY$_3$} is a pair $(X, N)$ with $X$ a
compact CY$_3$ and $N$ one of the $24$ Niemeier lattices (the Leech
lattice $\Lambda_{24}$ and the $23$ root-generated Niemeier lattices;
Conway--Sloane 1988). The Niemeier-lattice data $N$ determines:
\begin{itemize}
\item an arithmetic group $\Gamma_N \subset O(N \otimes \R)$ with
  $N$ as its lattice-polarisation;
\item a Borcherds automorphic form $\Phi_N$ of weight $\mathrm{rk}(N)/2$
  obtained by singular theta-lift from the weight-$1/2$ Vector-Valued
  modular form on $N$ (Scheithauer \texttt{arXiv:math/0409398}
  Thm.~3.1);
\item a BKM Lie algebra $\mathfrak{g}_N$ of Kac type governed by $N$'s
  root system.
\end{itemize}
\end{definition}

\begin{proposition}[$24$ BKM algebras at Niemeier cells $15$]
\label{prop:R13-niemeier-24-bkm}
\ClaimStatusProposition.
For each of the $24$ Niemeier lattices $N$ there is a distinct
BKM Lie algebra $\mathfrak{g}_N$ obtained by Borcherds lift of the
$N$-theta form; the corresponding cell-$15$ cocycle class
$[K^{(15)}_N]$ on a Niemeier-decorated CY$_3$ has $\kBKM$-invariant
\[
  \kBKM(\mathfrak{g}_N) \;=\; \frac{c_N^{\mathrm{Niemeier}}(0)}{2}
\]
where $c_N^{\mathrm{Niemeier}}(0)$ is the zeroth Fourier coefficient of
$\Phi_N$'s theta lift. Among the $24$:
\begin{itemize}
\item \emph{Leech $\Lambda_{24}$:} $\mathfrak{g}_{\Lambda_{24}}$ is the
  Fake Monster Lie algebra (Borcherds 1990 \emph{Nothing New}
  Chap.~6); $\kBKM = 12$ from $c(0) = 24$.
\item \emph{Niemeier $A_1^{24}$:} $\mathfrak{g}_{A_1^{24}}$ is a rank-$24$
  BKM; $\kBKM = 8$ via Scheithauer.
\item \emph{Niemeier $E_8^3$:} $\mathfrak{g}_{E_8^3}$ is a rank-$24$
  exceptional BKM; $\kBKM = 4$.
\item \emph{Niemeier $D_{24}$:} $\mathfrak{g}_{D_{24}}$ related to
  the spin-$1/2$ heterotic string, $\kBKM = 6$.
\item other $20$ Niemeiers: each contributes its own $\kBKM$-weight.
\end{itemize}
\end{proposition}

\begin{proof}
Scheithauer \texttt{arXiv:math/0409398}~\S 3 classifies Borcherds lifts
of reflective vector-valued modular forms on Niemeier lattices; each of
the $24$ Niemeiers supplies a unique BKM $\mathfrak{g}_N$ of rank equal
to $\mathrm{rk}(N) = 24$. The Fourier coefficient $c_N^{\mathrm{Niemeier}}(0)$
reads off the $\kBKM$-invariant via the Borcherds weight theorem
(CLAUDE.md: $\kBKM(\Phi_N) = c_N(0)/2$ universal). The Leech
entry yields the Fake Monster: Borcherds 1990 established
the Fake Monster as the Borcherds lift of the weight-$12$ theta
function of $\mathrm{II}_{25,1}$; the relevant CY dimension is $d = 5$
in Vol III's dimension-stratified picture (memory wave-15), not $d = 3$,
but the cell-$15$ cocycle classification on Niemeier-decorated
CY$_3$ records the Leech entry as a $24$-rank stratum-$(iv)$
cocycle fragment nonetheless (separately from the $d = 5$ envelope).
\end{proof}

\begin{theorem}[$24$-fold uniform cell-$15$ classification]
\label{thm:R13-24-fold-uniform-cell-15}
\ClaimStatusConjectured.
Let $X$ be a compact CY$_3$ admitting a Niemeier decoration $N$
(Definition~\ref{def:R13-niemeier-decorated-cy3}). The cell-$15$
cocycle class $[K^{(15)}]$ of $X$ is uniformly classified by $N$:
\[
  [K^{(15)}]_{(X, N)} \;\leftrightarrow\;
  [K^{(15)}_N] \;=\; (c^{(1)}_{24}, c^{(14)}_N) \;\in\;
  H^1(\cN_X, \underline{F}^{(15)}_N),
\]
with the Niemeier-dependent cocycle group
\[
  \underline{F}^{(15)}_N \;=\;
  \underline{\Pic}_{\Str_{i, \mathrm{loc}}} \rtimes
  \underline{\Aut_s(X)} \rtimes \underline{G_N} \rtimes
  \underline{\Gamma_N},
\]
in which $G_N$ is the symmetry group of $N$ (Conway group $\mathrm{Co}_0$
for Leech; $\mathrm{Aut}(N)$ for other Niemeier lattices) and
$\Gamma_N = O(N \otimes \R)$-stabiliser of the lattice. On
$X = K3 \times E$ with $N = \mathrm{II}_{2,19}$ the transcendental
lattice, the data reduces to the $\Delta_5$-Borcherds lift case of
Cycle~$4$ with $\kBKM = 5$.
\end{theorem}

\begin{proof}
The proof is assembled from:
\begin{itemize}
\item the K\"{u}nneth split of Theorem~\ref{thm:R13-kunneth-split},
  which applies to any $X = Y \times E$ with $Y$ a CY$_2$;
\item the $24$-germ toric stratum of Proposition~\ref{prop:R13-24-germ-stratum},
  generalised from $K3 \times E$ to any CY$_3$ with $N$-Niemeier
  decoration -- each of the $24$ Niemeier $N$'s corresponds to a distinct
  $(24-\text{generic-fibre-stalk})$ decomposition of the elliptic/K3-type
  fibration;
\item the Chen--Ruan Mathieu closure of Theorem~\ref{thm:R13-mathieu-dbar-closed},
  generalised to the $N$-Niemeier orbifold by replacing $M_{24}$ by
  the appropriate $G_N$ symmetry group;
\item the Borcherds--Heegner cocycle of
  Proposition~\ref{prop:R13-borcherds-heegner}, generalised to the
  $\Phi_N$-Borcherds lift on the Niemeier period domain.
\end{itemize}
The Niemeier classification provides the unified framework for $24$
distinct CY$_3$-witnesses, all at cell $15$, each with its own
$\kBKM$-weight $c_N(0)/2$.
\end{proof}

\begin{remark}[Dimensional-stratification cross-check]
\label{rem:R13-dim-strat-crosscheck}
The Leech $\Lambda_{24}$-Niemeier, whose BKM is the Fake Monster,
lives canonically at $d = 5$ in the Vol~III dimensional stratification
(wave-15 memory): rank $25 > 20 = h^{1,1}(K3)$, so no $d = 3$
embedding exists; Dunn--Lurie $E_5 \simeq E_2 \otimes E_2 \otimes E_1$
places the Fake Monster on $K3 \times K3 \times E$ at $d = 5$. The
Niemeier-classification theorem above records the Leech/Fake-Monster
entry as a $d = 5$ sibling of the $d = 3$ entries; the cell-$15$
cocycle structure nonetheless includes the Leech as a uniformly
classified entry, with the caveat that the physical-CY$_3$ witness
for the Fake Monster does not exist (only the $d = 5$ one does).
The remaining $23$ root-generated Niemeiers give genuine $d = 3$
CY$_3$-witnesses when tensored with an elliptic curve and restricted
to appropriate Heegner-divisor loci.
\end{remark}

\subsection*{Surviving core of (A5)}

The cell-$15$ classification extends uniformly across the $24$ Niemeier
lattices (Leech plus $23$ root-generated), giving $24$ distinct
CY$_3$-witnesses with inequivalent $\kBKM$-invariants read off from
the zeroth Fourier coefficient of each Borcherds lift $\Phi_N$. The
$K3 \times E$-witness of Cycle~$4$ is the $\mathrm{II}_{2,19}$-entry
with $\kBKM = 5$; the Leech-entry $\kBKM = 12$ corresponds to the
Fake Monster at $d = 5$, a cross-dimensional sibling.

% =====================================================================
\section{Cross-volume corridor, four mirror statements}
% =====================================================================

The explicit cell-$14$/$15$ cocycle computation locks into the
three-volume programme as follows.

\paragraph{Vol I (bar--cobar).} Under Stage~2 specialisation
$\mathrm{Sp}_{\Sigma_2, C}$, the cell-$15$ cocycle pushes forward to a
bar-cobar cocycle on the $E_1$-chiral shadow; the $48$-rank local-toric-germ
contribution becomes a curve-stalk $1$-cocycle on the reference curve
$C$, one curve-stalk per $I_1$-fibre image under the Ran projection.

\paragraph{Vol II (3D HT QFT).} The Vol~II $E_2$-hFA on $\Sigma_2 \times C$
factors through Dunn $E_3 \simeq E_2 \otimes E_1$; restricting the
cell-$15$ cocycle to $\Sigma_2 \times C \subset K3 \times E$ lands in
the Vol~II geometric engine, transverse to $\Sigma_2$. The Mathieu
Chen--Ruan twist becomes an equivariant surface-integral contribution
to the Vol~II construction.

\paragraph{Vol III (CY quantum groups).} The seven $r_{\mathrm{CY}}$-face
calculations on $K3 \times E$ each read off some projection of the
cell-$15$ cocycle; the $\kBKM = 5$ value is the BKM-axis face, read off
from the Borcherds singular theta-lift. The super-Yangian
$Y_\hbar(\fso(4, 20))$ of Proposition~\ref{gluing:prop:super-yangian-emergence}
is comparison-target read off from the Hodge-involution component of
the cell-$15$ cocycle.

\paragraph{Pattern 273.} The object-level cell-$15$ cocycle is the
explicit \v{C}ech $1$-cocycle on the product nerve
$\cN_{K3}^{\mathrm{CR}} \times \cN_E$; the functorial cell-$15$
classification assigns, to each compact Niemeier-decorated CY$_3$,
its uniform cell-$15$ cocycle class. The two are Pattern-273 siblings.

% =====================================================================
\section{Status summary}
% =====================================================================

\begin{enumerate}[leftmargin=2em]
\item Theorem~\ref{thm:R13-kunneth-split} supplies the K\"{u}nneth split
of the cell-$14$/$15$ transition FM-kernel on the product-compatible
cover of $K3 \times E$.
\item Corollary~\ref{cor:R13-cell-14-rep} gives an explicit \v{C}ech
$1$-cocycle representative for cell $14$; Corollary~\ref{cor:R13-cell-15-rep}
adds the $24$-local-toric-germ contribution for cell $15$.
\item Theorem~\ref{thm:R13-mathieu-dbar-closed} resolves the
mock-modular anomaly of the Mathieu twining by passage to Chen--Ruan
orbifold inertia, supplying $\bar\partial$-closedness of the cell-$14$
cocycle on the inertia-refined nerve.
\item Theorem~\ref{thm:R13-leray-ss} gives the Leray spectral sequence
computation on the product nerve; degenerates at $E_2$ for the split
sheaf.
\item Proposition~\ref{prop:R13-24-germ-stratum} and Corollary~\ref{cor:R13-toric-input-E2-00}
establish the $48$-rank local-toric-germ contribution at cell $15$.
\item Proposition~\ref{prop:R13-borcherds-heegner} and
Theorem~\ref{thm:R13-vanishing-verdict} give the vanishing/non-vanishing
verdict: cell $14$ and cell $15$ are both non-zero on generic
$K3 \times E$; the paramodular-quotient absorbs all Hecke/Mathieu
freedom to the single $\kBKM(\mathfrak{g}_{\Delta_5}) = 5$ class.
\item Theorem~\ref{thm:R13-24-fold-uniform-cell-15} extends the cell-$15$
classification uniformly across the $24$ Niemeier lattices, giving
$24$ distinct CY$_3$-witnesses with inequivalent $\kBKM$-invariants.
\end{enumerate}

The construction is compatible with all prior Vol~III content
(fifteen-cell classification, product-compatible cover,
Serre-equivariant quasi-NCCR, Borcherds weight universality, Mathieu
moonshine) and supplies the explicit \v{C}ech cocycle representatives
for the deepest cells of the stratum-occupancy classification.

\section*{Four frontier obstructions beyond the scope}

\paragraph{(R13-F1) Explicit computation of the $H^1(\cN_{K3}, \underline{\Aut_s})$
of the Mathieu K3.} The cell-$14$ $K3$-factor cocycle
$c^{K3, \mathrm{twin}}_{\alpha \alpha'}$ is given abstractly as an
Eguchi--Ooguri--Tachikawa mock-modular cocycle; the explicit Euler-characteristic
computation of $H^1(\cN_{K3}, \underline{\Aut_s})$ (dimension,
generator-relation presentation) requires a fine Mathieu-symmetric
analysis of the $\cN_{K3}$-nerve structure.

\paragraph{(R13-F2) Explicit Atkin--Lehner cocycle on $\cN_E$.} The
cell-$14$ $E$-factor cocycle $c^{E, \mathrm{Hecke}}_{\beta \beta'}$ is
given as an Atkin--Lehner Hecke cocycle; for $E$ non-CM it is zero,
but for CM-$E$ it is non-zero with values in the Hecke-monodromy
group. The explicit dependence on the CM-discriminant remains a
computational task.

\paragraph{(R13-F3) Non-abelian $M_{24}$-inertia stratum.} The
Chen--Ruan inertia refinement $I([K3/G])$ of Cycle~$2$ assumes
$G$ abelian or a well-behaved non-abelian subgroup of $M_{24}$;
for $G \subset M_{24}$ non-abelian (e.g.\ $M_{23}$, $M_{22}$), the
inertia stack carries non-trivial $2$-form information from the
non-abelian conjugacy structure, which the current construction does
not fully track.

\paragraph{(R13-F4) $23$-Niemeier witness landscape.} Theorem~\ref{thm:R13-24-fold-uniform-cell-15}
asserts a uniform cell-$15$ classification across the $24$ Niemeier
lattices. For explicit CY$_3$-witnesses at each Niemeier $N \neq
\mathrm{II}_{2,19}$, specific compact CY$_3$ examples remain to be
identified. The landscape of CY$_3$'s carrying $N$-Niemeier
decorations is not exhausted by $K3 \times E$.

\section*{Epilogue: deep-semantic-merge with Agent 14}

Agent 14 established the abstract fifteen-cell classification of
transition-FM-kernel cocycles for holomorphic $E_3$-factorisation
data on compact CY$_3$, with cells $14$ and $15$ identified as the
deepest-strata-occupancy cells on $K3 \times E$. The present wave
supplies the missing explicit content: \v{C}ech $1$-cocycle
representatives for cells $14$ and $15$, Leray spectral sequence
computation of the classifying cohomology, Chen--Ruan inertia
refinement to restore $\bar\partial$-closedness, $24$-local-toric-germ
contribution from the $I_1$-fibre stratum, vanishing verdict showing
both cells non-zero on generic $K3 \times E$, and extension to the
Niemeier-uniform classification across $24$ Niemeier lattices.

The present content is strictly stronger than Agent 14's stratum-occupancy
table: Agent 14 named the shape of the cocycle group, this wave
supplies the explicit cocycle, its cohomology-group computation, its
vanishing behaviour, and its Niemeier-uniform extension. The
local-to-global descent of $\PhiFA_3(\Perf(X))$ on compact CY$_3$ is
thus instantiated, for the $K3 \times E$-class, by an explicit
\v{C}ech cocycle with computable cohomology and verifiable
non-vanishing.

\bigskip

\noindent \emph{End of Agent R-13, wave residuals.}

\end{document}
